\begin{trapbox}
\begin{center}
\begin{tikzpicture}[font=\small]
 \path[use as bounding box] (-3,-3) rectangle (9,3);
 \foreach \sx/\endpoint in {0/2,6/1}{
 \begin{scope}[shift={(\sx,0)}]
 \draw[blue!60!black,thick] (0,0) ellipse (2 and {sqrt(3)});
 \draw[gray] (-2.4,0)--(2.4,0);
 \draw[teal!70!black,thick] (-\endpoint,0)--(0,{sqrt(3)});
 \draw[magenta!70!black,thick] (\endpoint,0)--(0,{sqrt(3)});
 \fill (-\endpoint,0) circle (1.7pt);\fill (\endpoint,0) circle (1.7pt);
 \node[above] at (0,{sqrt(3)}) {$P(0,\sqrt3)$};
 \end{scope}}
 \node at (0,2.65) {指定长轴顶点};\node at (6,2.65) {换成焦点：定值不适用};
 \node[below left] at (-2,0) {$A_1$};\node[below right] at (2,0) {$A_2$};
 \node[below] at (5,0) {$F_1$};\node[below] at (7,0) {$F_2$};
 \node at (0,-2.3) {$k_{PA_1}k_{PA_2}=-\frac34$};
 \node at (6,-2.3) {$k_{PF_1}k_{PF_2}=-3$};
 \node at (3,-2.9) {同一椭圆 $x^2/4+y^2/3=1$；两图单位长度相同};
\end{tikzpicture}
\end{center}


\begin{description}[style=nextline,leftmargin=2.8em,labelsep=0.5em,itemsep=0.4em]
 \item[\textbf{\textcolor{CTrap}{陷阱 1：忽视斜率存在的范围}}]
  \textbf{错误：}以为 $P$ 在椭圆上就可无条件套式，包含 $P=A_1,A_2$。\\
  \textbf{纠正：}端点处某一条连线退化，分母为零，不能写“斜率积”。
  反向求轨迹时也要从椭圆中删掉这两点。
 \item[\textbf{\textcolor{CTrap}{陷阱 2：乘积的负号丢失}}]
  \textbf{错误：}$\frac{y^2}{x^2-a^2}=+\frac{b^2}{a^2}$。\\
  \textbf{纠正：}$a^2-x^2=-(x^2-a^2)$，两斜率一正一负，结果必为负。
 \item[\textbf{\textcolor{CTrap}{陷阱 3：横轴与纵轴公式混用}}]
  \textbf{错误：}竖轴椭圆仍套 $-b^2/a^2$。\\
  \textbf{纠正：}竖轴端点 $B_{1,2}=(0,\mp a)$，斜率写为 $(y\pm a)/x$，
  所以结果是 $-a^2/b^2$，不是 $-b^2/a^2$。
 \item[\textbf{\textcolor{CTrap}{陷阱 4：误把“长轴顶点”当成“焦点”}}]
  \textbf{错误：}把端点取为 $F_1(-c,0),F_2(c,0)$，仍写相同定值。\\
  \textbf{纠正：}本结论所用两点为 $A_1(-a,0),A_2(a,0)$，不是 $F_1,F_2$；
  任意其他两个定点一般不满足同样的斜率积公式。
 \item[\textbf{\textcolor{CTrap}{陷阱 5：只知斜率积，忘查点是否在椭圆上}}]
  \textbf{最小反例：}对 $\frac{x^2}{4}+\frac{y^2}{3}=1$，取不在椭圆上的 $P=(0,1)$，
  斜率积为 $(\frac12)(-\frac12)=-\frac14$，而目标值是 $-\frac34$。
 \item[\textbf{\textcolor{CTrap}{陷阱 6：逆命题忘查长短轴}}]
  若 $k_{PA_1}k_{PA_2}=-\lambda$，轨迹为
  $\frac{x^2}{a^2}+\frac{y^2}{\lambda a^2}=1$（删去端点）。
  $\lambda>1$ 时，$A_1A_2$ 对应的是\textbf{短轴}，不能再把这两个点叫“长轴端点”。
 \item[\textbf{\textcolor{CTrap}{陷阱 7：以为坐标轴可以任意旋转}}]
  \textbf{纠正：}本式使用的是常见 $y$ 对 $x$ 的斜率；
  椭圆转到倾斜位置后不能机械沿用“$-b^2/a^2$”，须先考虑坐标系或重新计算。
 \item[\textbf{\textcolor{CTrap}{陷阱 8：把“第三定义”误认为唯一正式定义}}]
  这里“第三定义”是教材外常用的\textbf{二级结论名称}，
  正式题目中应交代椭圆标准方程、端点与点的位置，不宜只写名称而不说明条件。
 \item[\textbf{\textcolor{CTrap}{陷阱 9：认为所有圆点都能“随便”配到椭圆点}}]
  \textbf{错误：}只说“圆变椭圆”，便把圆上 $Q(u,v)$ 配到椭圆上的任意点 $P$。\\
  \textbf{纠正：}要先说明\textbf{选用了哪一种映射}。
  本证明特选“横坐标不变、上下半平面不互换”，再由椭圆方程推出
  $P=(u,\frac ba v)$。这种映射方便固定 $A_1,A_2$，但不是唯一可能的映射。
 \item[\textbf{\textcolor{CTrap}{陷阱 10：把纵向压缩当作等比缩放，认为直角还在}}]
  \textbf{错误：}由 $QA_1\perp QA_2$ 直接推出 $PA_1\perp PA_2$。\\
  \textbf{纠正：}只在纵方向乘 $b/a$，与横方向的比例不同，\textbf{角度一般不保留}。
  $PA_1,PA_2$ 的斜率积是 $-(b/a)^2$，一般不等于 $-1$。
 \item[\textbf{\textcolor{CTrap}{陷阱 11：弄反映射或不注意上下符号}}]
  \textbf{错误：}圆 $\to$ 椭圆时用 $a/b$，或把上半圆的点映射到下半椭圆。\\
  \textbf{纠正：}$T(u,v)=(u,\frac ba v)$ 是压缩；
  $T^{-1}(x,y)=(x,\frac ab y)$ 是拉伸。所选的正比例 $b/a$ 保持 $v$ 与 $y$ 同号；
  $\frac ya=\frac vb$ 不是本映射的正确归一化关系，正确的是 $\boxed{\frac yb=\frac va}$。
\end{description}
\end{trapbox}
